% === Preamble for Engineering ===

% --- Encoding and fonts ---
\usepackage[utf8]{inputenc}
\usepackage[T1]{fontenc}
\usepackage{lmodern}

% --- Mathematics ---
\usepackage{amsmath}
\usepackage{amssymb}
\usepackage{amsthm}
\usepackage{mathtools}
\usepackage{bm}
\usepackage{mathrsfs}
\usepackage{stmaryrd}
\usepackage{siunitx}
\sisetup{detect-all = true}

% --- Graphics and diagrams ---
\usepackage{tikz}
\usetikzlibrary{
  arrows.meta,
  calc,
  positioning,
  shapes.geometric,
  shapes.misc,
  fit,
  backgrounds,
  decorations.markings,
  decorations.pathreplacing,
  matrix,
  chains,
  graphs,
  quotes,
  circuits.ee.IEC,
}
\usepackage{pgfplots}
\pgfplotsset{compat=1.18}

% --- Tables ---
\usepackage{booktabs}
\usepackage{array}
\usepackage{longtable}

% --- Code and pseudocode ---
\usepackage{algorithm}
\usepackage{algpseudocode}

% --- Colors ---
\usepackage{xcolor}
\definecolor{engblack}{RGB}{0, 0, 0}
\definecolor{englime}{RGB}{132, 204, 22}     % #84cc16, the institute accent
\definecolor{engblue}{RGB}{31, 78, 121}
\definecolor{engred}{RGB}{153, 42, 42}
\definecolor{engamber}{RGB}{180, 110, 30}
\definecolor{enggray}{RGB}{88, 88, 88}

% --- Bibliography ---
\usepackage[
  backend=biber,
  style=authoryear,
  sorting=nyt,
  maxbibnames=99,
  maxcitenames=2,
  giveninits=true,
  uniquename=init,
  dashed=false,
]{biblatex}
\addbibresource{bibliography/references.bib}

% --- Cross-references and links ---
\usepackage{hyperref}
\hypersetup{
  colorlinks=true,
  linkcolor=engblue,
  citecolor=engamber,
  urlcolor=engred,
}
\usepackage[capitalize,noabbrev]{cleveref}

% Repository-relative paths in prose (e.g., docs/research/accidents/...).
% \path{} from the url package (loaded by hyperref) breaks at slashes,
% hyphens, and dots, preventing long-path overflow. We wrap it as
% \repopath{} so the semantic intent (repository path, not URL) is
% explicit at the call site.
\newcommand{\repopath}[1]{\path{#1}}

% --- Indexing ---
\usepackage{makeidx}
\makeindex

% --- Table of contents spacing ---
% Part numbers reach four characters ("VIII"), chapter numbers run globally
% across all twelve books and reach three digits (max 163), and section
% labels reach six characters ("163.10"). The default memoir indent widths
% are sized for shorter labels and collide. Widen the number columns at
% every depth.
\cftsetindents{part}{0em}{3.0em}
\cftsetindents{chapter}{0em}{3.2em}
\cftsetindents{section}{3.2em}{4.2em}
\cftsetindents{subsection}{7.4em}{5.0em}

% --- Epigraphs ---
\usepackage{epigraph}
\setlength{\epigraphwidth}{0.7\textwidth}

% --- Miscellaneous ---
\usepackage[inline]{enumitem}
\usepackage{microtype}

% === Theorem-like environments ===
\theoremstyle{plain}
\newtheorem{theorem}{Theorem}[chapter]
\newtheorem{proposition}[theorem]{Proposition}
\newtheorem{lemma}[theorem]{Lemma}
\newtheorem{corollary}[theorem]{Corollary}
\newtheorem{conjecture}[theorem]{Conjecture}

\theoremstyle{definition}
\newtheorem{definition}[theorem]{Definition}
\newtheorem{axiom}[theorem]{Axiom}
\newtheorem{principle}[theorem]{Principle}
\newtheorem{archetype}[theorem]{Archetype}
\newtheorem{model}[theorem]{Model}

\theoremstyle{remark}
\newtheorem{remark}[theorem]{Remark}
\newtheorem{example}[theorem]{Example}
\newtheorem{exercise}{Exercise}[chapter]
\newtheorem{project}{Project}[chapter]
\newtheorem{lab}{Lab}[chapter]
\newtheorem{estimation}{Estimation}[chapter]   % Q20: every chapter has these.

% === Chapter metadata block ===
% Renders a small italic gray paragraph immediately under the chapter title
% with the chapter's editorial metadata: half-life category, archetypes
% invoked, project track (physical/simulation/household/hybrid), exercise
% count target, and named accident cases (where applicable).
%
% Usage in chapter source:
%   \chapmeta{Half-life: foundational. Archetypes: scaling, failure.
%            Project track: physical. Exercise target: 28.
%            Named cases: Mars Climate Orbiter, Air Canada 143, Ariane 5.}
%
% The block is visible to the writer and reviewer during drafting. A future
% release pass may hide it via a class option; for now keep it visible.
\newcommand{\chapmeta}[1]{%
  \begingroup
  \par\medskip
  \footnotesize\itshape\color{enggray}%
  #1%
  \par\endgroup
  \medskip
}

% === Reader-path tagging (settled Q51, 2026-04-28) ===
% Each section is one of three reader paths:
%   core      = the spine; required to call the book read.
%   standard  = working-engineer depth; the chapter's actual argument.
%   mastery   = optional advanced material; reader may skip without losing the thread.
%
% Usage:
%   \section{What measurement is for}\pathtag{core}
%   \begin{mastery} ... optional inline advanced box ... \end{mastery}
\newcommand{\pathtag}[1]{%
  \par\nobreak
  {\footnotesize\itshape\color{enggray}\textsf{[Path: #1]}}\par
  \nobreak\smallskip
}

\newenvironment{mastery}{%
  \par\medskip
  \noindent\rule{\textwidth}{0.4pt}\par\nobreak
  \begingroup\small\color{enggray}%
  \noindent\textsf{\textbf{Mastery (optional).}}\par\nobreak\smallskip
}{%
  \par\endgroup
  \nobreak\noindent\rule{\textwidth}{0.4pt}\par\medskip
}

% === Part interludes ===
\newenvironment{interlude}{%
  \cleardoublepage
  \thispagestyle{empty}
  \vspace*{3cm}
  \begin{center}\rule{4cm}{0.4pt}\end{center}
  \vspace{0.5cm}
  \begin{quote}\itshape
}{%
  \end{quote}
  \vspace{1cm}
  \begin{center}\rule{4cm}{0.4pt}\end{center}
  \clearpage
}

% === Cross-domain notation ===
\usepackage{eng-macros}

% === Memoir adjustments ===

% Custom chapter style that handles three-digit chapter numbers cleanly.
% Chapters are numbered globally across all twelve books (1..163 in the current
% structure), so the chapter heading must give the number its own line and a
% generous gap to the title. The default memoir styles (veelo, default, thatcher)
% either overlap the number with the title at three digits or push the number
% into the page margin.
\makechapterstyle{engineering}{%
  \renewcommand*{\chapnumfont}{\normalfont\HUGE\bfseries}
  \renewcommand*{\chaptitlefont}{\normalfont\Huge\bfseries\raggedright}
  \setlength{\beforechapskip}{60pt}
  \setlength{\midchapskip}{18pt}
  \setlength{\afterchapskip}{40pt}
  \renewcommand*{\printchaptername}{}
  \renewcommand*{\chapternamenum}{}
  \renewcommand*{\printchapternum}{%
    {\chapnumfont\thechapter}\par\nobreak\vskip\midchapskip}
  \renewcommand*{\afterchapternum}{}
  \renewcommand*{\printchaptertitle}[1]{%
    \hrule\vskip8pt\chaptitlefont ##1}%
}

\chapterstyle{engineering}
\pagestyle{headings}

% The twelve top-level divisions are called Volumes, not Parts. Engineering
% is the book; the twelve are its volumes.
\renewcommand{\partname}{Volume}
\renewcommand{\partnamefont}{\normalfont\Large\scshape}
\renewcommand{\partnumfont}{\normalfont\Large\scshape}
\renewcommand{\parttitlefont}{\normalfont\Huge\bfseries}

\setsecnumdepth{subsection}
\settocdepth{subsection}
